% Subsection: Arbitrary Operations  (notes stub; content authored later)
\subsection{Arbitrary Operations}

\begin{remark*}[Exposition]
Arbitrary operations allow an indexed family of inputs with no finite or
countable restriction on the indexing set.  They are broader than countable
operations and therefore impose stronger demands when a set system is required
to be closed under them.
\end{remark*}

\begin{definition}[Arbitrary Operation on a Family]\label{def:arbitrary-operation-on-family}
Let \(X\) be a set and let \(\mathcal{F}\subseteq\mathcal{P}(X)\).  An
\emph{arbitrary operation on \(\mathcal{F}\)} is an operation whose inputs are
an indexed family \(\{A_\alpha\}_{\alpha\in I}\) of members of
\(\mathcal{F}\), where \(I\) is any indexing set, and whose output is a subset
of \(X\).
\end{definition}

\begin{remark*}[Standard quantified statement]
The standard arbitrary union and arbitrary intersection operations have the
forms
\[
  \{A_\alpha\}_{\alpha\in I}\subseteq\mathcal{F}
  \quad\Longmapsto\quad
  \bigcup_{\alpha\in I}A_\alpha\subseteq X
\]
and
\[
  \{A_\alpha\}_{\alpha\in I}\subseteq\mathcal{F}
  \quad\Longmapsto\quad
  \bigcap_{\alpha\in I}A_\alpha\subseteq X.
\]
\end{remark*}

\begin{remark*}[Interpretation]
Every finite or countable operation is an arbitrary operation for a suitable
indexing set, but not conversely.  Arbitrary unions are central in topology.
Arbitrary intersections are also the mechanism behind generated structures:
one often intersects all admissible systems satisfying a condition in order to
obtain the smallest such system.
\end{remark*}

\begin{remark*}[Non-Examples]
A countable union is not arbitrary because it uses a countable index set; it is
arbitrary only in the broader sense that every countable index set is still an
index set.  The distinction matters when a definition permits arbitrary
unions but does not permit arbitrary intersections, or permits countable
unions but not arbitrary unions.
\end{remark*}

\begin{dependencies}
\begin{itemize}
  \item \hyperref[def:family-of-subsets]{Family of Subsets}
  \item \hyperref[def:indexed-family]{Indexed Family of Sets}
  \item \hyperref[def:indexed-union]{Indexed Union}
  \item \hyperref[def:indexed-intersection]{Indexed Intersection}
\end{itemize}
\end{dependencies}
